
\section{Methodology}

\noindent This section describes the data foundation, processing pipeline, and implementation choices used to transform clinical consultation audio into validated HL7 FHIR resources. It first motivates the selected data source and its preparation, then outlines the six-stage pipeline, and finally details the design decisions and runtime environment that support end-to-end processing.

\subsection{Data Source and Preparation}

PriMock57 \cite{papadopoulos2022primock57} was selected because it provides a reproducible and clinically relevant benchmark for evaluating the proposed audio-to-FHIR pipeline. Its paired audio and transcript annotations make it suitable for assessing how transcription, post-processing, and structured extraction interact across the full workflow rather than in isolation.

The PriMock57 dataset consists of 57 mock primary care consultations recorded over five days by clinicians at Babylon Health. The corpus comprises 114 mono WAV audio files: a separate doctor track and a separate patient track per consultation, named according to the convention \path{dayN_consultationNN_role.wav}, where \textit{N} indexes the recording day (1–5) and \textit{role} is \texttt{doctor} or \texttt{patient}. Ground-truth annotations include manual utterance-level transcriptions and consultation notes produced by the recording clinicians. PriMock57 is, to our knowledge, the only publicly available clinical audio benchmark with ground-truth transcripts at the consultation level, making it the natural evaluation corpus for end-to-end audio-to-FHIR research.

Prior to pipeline ingestion, the two per-consultation tracks are mixed via additive synthesis: the doctor and patient waveforms are loaded at the same sample rate, summed sample-by-sample, and peak-normalized to amplitude 1.0. This yields a single-channel recording that simulates a room-microphone capture of the consultation dyad. An optional 60-second trim was applied to individual consultations during development to accelerate iterative testing.

\subsection{Pipeline Architecture}

The system is organized as a six-step pipeline spanning four functional phases: (1) acoustic pre-processing, (2) automatic speech recognition with speaker attribution, (3) LLM-driven transcript refinement and clinical information extraction, and (4) FHIR resource generation with validation and provenance tracking. Figure~\ref{fig:pipeline} presents this design as a workflow, emphasizing how audio and intermediate artifacts move through the pipeline from raw consultation recording to validated structured output.

\begin{enumerate}
  \item Audio Ingestion \& Pre-processing
  \item Speaker Diarization
  \item ASR Transcription
  \item LLM Transcript Cleaning \& Speaker Role Assignment
  \item LLM Clinical Extraction \textrightarrow{} FHIR R4 Bundle Construction
  \item FHIR Validation \& Provenance Tracking
\end{enumerate}

Figure~\ref{fig:pipeline} summarizes how these stages connect in sequence. The workflow begins with consultation audio ingestion and pre-processing, followed by speaker diarization, segment-level transcription, and transcript assembly. The resulting transcript is then refined with LLM-based cleaning and speaker role assignment before a second LLM pass extracts clinical content and assembles a FHIR R4 Bundle through programmatic \texttt{fhir.resources} construction. The final stage validates the generated resources and appends provenance metadata so each output remains auditable.

\begin{figure}[ht]
\centering
\resizebox{\linewidth}{!}{%
\begin{tikzpicture}[
    node distance=1.35cm and 1.35cm,
    process/.style={
        rectangle, rounded corners=7pt,
        minimum width=4.0cm, minimum height=1.35cm,
        text centered, text width=3.7cm,
        font=\normalsize\bfseries,
        line width=1.2pt
    },
    artifact/.style={
        trapezium, trapezium left angle=75, trapezium right angle=105,
        minimum width=3.8cm, minimum height=1.15cm,
        text centered, text width=3.3cm,
        font=\normalsize\bfseries,
        line width=1.1pt
    },
    terminal/.style={
        ellipse,
        minimum width=3.7cm, minimum height=1.1cm,
        text centered,
        font=\normalsize\bfseries,
        line width=1.2pt
    },
    audio/.style   ={process, fill=blue!15,   draw=blue!60!black},
    asr/.style     ={process, fill=cyan!20,   draw=cyan!60!black},
    llm/.style     ={process, fill=teal!20,   draw=teal!60!black},
    fhir/.style    ={process, fill=green!18,  draw=green!50!black},
    data/.style    ={artifact, fill=gray!10,  draw=gray!60!black},
    terminus/.style={terminal, fill=orange!18, draw=orange!65!black},
    note/.style={
        font=\small\itshape,
        text=gray!75!black,
        text centered
    },
    flow/.style={
        -{Stealth[length=8pt,width=6pt]},
        line width=1.6pt,
        color=gray!70!black
    },
]

\node[terminus] (input) {Consultation Audio Input};
\node[note, below=0.08cm of input] {mixed doctor/patient recording};

\node[audio, right=1.6cm of input] (step1) {Step 1: Audio Ingestion \& Pre-processing};
\node[note, below=0.08cm of step1] {16\,kHz mono, denoising, VAD};

\node[asr, right=1.6cm of step1] (step2) {Step 2: Speaker Diarization};
\node[note, below=0.08cm of step2] {pyannote.audio, max 2 speakers};

\node[data, below=1.7cm of step2] (segments) {Diarized Speech Segments};
\node[note, below=0.08cm of segments] {timestamped speaker-labeled intervals};

\node[asr, left=1.6cm of segments] (step3) {Step 3: ASR Transcription};
\node[note, below=0.08cm of step3] {Medical-Whisper-Large-v3};

\node[data, left=1.6cm of step3] (transcript) {Raw Transcript Segments};
\node[note, below=0.08cm of transcript] {text with generic speaker IDs};

\node[llm, below=1.8cm of transcript] (step4) {Step 4: Transcript Refinement \& Role Assignment};
\node[note, below=0.08cm of step4] {cleaning and physician/patient mapping};

\node[data, right=1.6cm of step4] (cleaned) {Post-processed Transcript};
\node[note, below=0.08cm of cleaned] {clean text plus source traceability};

\node[llm, right=1.6cm of cleaned] (step5) {Step 5: Clinical Extraction \& FHIR Construction};
\node[note, below=0.08cm of step5] {LLM extraction with \texttt{fhir.resources}};

\node[fhir, right=1.6cm of step5] (step6) {Step 6: Validation \& Provenance};
\node[note, below=0.08cm of step6] {R4B checks and provenance records};

\node[terminus, below=1.8cm of step5] (output) {Validated FHIR Bundle Output};
\node[note, below=0.08cm of output] {Bundle plus per-step JSON artifacts};

\draw[flow] (input) -- (step1);
\draw[flow] (step1) -- (step2);
\draw[flow] (step2) -- (segments);
\draw[flow] (segments) -- (step3);
\draw[flow] (step3) -- (transcript);
\draw[flow] (transcript) -- (step4);
\draw[flow] (step4) -- (cleaned);
\draw[flow] (cleaned) -- (step5);
\draw[flow] (step5) -- (step6);
\draw[flow] (step6) |- (output);

\begin{scope}[on background layer]
    \node[
        fill=cyan!6, draw=cyan!30, rounded corners=9pt, line width=0.9pt,
        fit=(step1)(step2)(segments)(step3)(transcript),
        inner sep=9pt,
        label={[font=\scriptsize\itshape, text=cyan!55!black]above:Speech Processing Workflow}
    ] {};

    \node[
        fill=orange!6, draw=orange!30, rounded corners=9pt, line width=0.9pt,
        fit=(step4)(cleaned)(step5)(step6)(output),
        inner sep=9pt,
        label={[font=\scriptsize\itshape, text=orange!55!black]below:LLM-to-FHIR Workflow}
    ] {};
\end{scope}

\end{tikzpicture}%
}% end resizebox

\caption{Workflow for the proposed six-step clinical audio-to-FHIR pipeline.}
\label{fig:pipeline}
\end{figure}

A central design decision is the collapse of the traditional multi-stage NLP pipeline (typically comprising SOAP segmentation, named entity recognition, negation detection, terminology normalization, and FHIR mapping as discrete components) into a single LLM extraction pass followed by schema-constrained, programmatic FHIR construction (Steps 5–6). This consolidation is motivated by Frei et al., who demonstrated a hallucination rate below 2.3\% when FHIR resources are assembled programmatically via \texttt{fhir.resources} object construction rather than generated as free-form JSON \cite{frei2025inferno}, and by Li et al., who demonstrated the feasibility of joint LLM-driven extraction and FHIR transformation \cite{li2024fhirgpt}.

All intermediate outputs are persisted to disk after each step (\path{outputs/<name>/step0N_<name>.json}). Completed steps are detected at startup and skipped, enabling iterative development and evaluation without full pipeline recomputation.

\subsubsection{Audio Ingestion and Pre-processing (Step 1)}

Step 1 accepts a single audio file and produces a pre-processed float32 waveform suitable for downstream ASR and diarization, using three libraries: librosa ($\geq$0.10.2) for loading, noisereduce ($\geq$3.0.0) for stationary noise suppression, and webrtcvad-wheels for voice activity detection (VAD). The mixed audio file is loaded at 16,000\,Hz and downmixed to mono via librosa. Stationary background noise is suppressed using noisereduce's spectral subtraction method, which estimates a noise profile from nominally silent regions. Voice activity is then detected using WebRTC VAD, applied over non-overlapping 20\,ms frames (320 samples at 16\,kHz) at aggressiveness level 2 on the 0--3 scale. Frames classified as voiced are retained; the remainder are discarded. The output is an \texttt{IngestionResult} containing the voiced float32 array, sample rate, total duration, \textit{speech\_ratio} (voiced frames / total frames), and source path.

Empirically, speech\_ratio $\approx 0.85$ was observed on \texttt{day1\_consultation01\_mixed.wav} (387\,s total; approximately 328\,s voiced content). Filtering silence at this stage reduces ASR transcription time and the number of non-speech segments passed to the diarization module.

\subsubsection{Speaker Diarization (Step 2)}

Speaker diarization is performed using pyannote.audio 4.x \cite{tran2022asr} with the \path{pyannote/speaker-diarization-community-1} pre-trained model, executed on CPU. The pipeline is constrained to a maximum of two speakers, reflecting the doctor–patient consultation dyad. Output is a list of \texttt{Segment(start, end, speaker)} objects with generic labels \texttt{SPEAKER\_00} and \texttt{SPEAKER\_01}. Speaker role assignment (mapping generic labels to \textsc{physician} or \textsc{patient}) is deferred to Step 4, where the full transcript context is available to enable more accurate role inference. Independent evaluation on clinical dialogue has reported word-level diarization error rates (WDER) of 1.8–13.9\% under controlled conditions \cite{tran2022asr}.

\subsubsection{ASR Transcription (Step 3)}

Transcription uses \path{Na0s/Medical-Whisper-Large-v3}, a 1.55B-parameter model fine-tuned on the PriMock57 training split by Banerjee et al.\ \cite{banerjee2024medasr}. The fine-tuned model achieves a word error rate of 0.19 on the PriMock57 validation split and 0.24 on the test split, representing a 42\% relative WER improvement over the zero-shot Whisper large-v3 baseline (WER 0.33) \cite{radford2023whisper}. The model is loaded via the HuggingFace \texttt{transformers} pipeline using \texttt{torch.float32} on CPU.

For each diarized segment of duration $\geq 0.1$\,s, the corresponding audio is sliced from the waveform using integer sample indices at 16\,kHz and passed to the ASR pipeline with \texttt{language="en"}. Each transcription is emitted as a \texttt{TranscriptSegment} (fields: start, end, speaker, text) carrying the diarization label. Segments shorter than 0.1\,s are discarded to avoid transcription artefacts.

\subsubsection{Transcript Cleaning and Speaker Role Assignment (Step 4)}

Step 4 applies GPT-5-mini via the agno Agent framework with Pydantic-enforced structured output. A single LLM pass over the complete transcript performs two tasks. First, transcription cleaning: filler words are removed, medical homophones are resolved (e.g.\ \textit{ileum} vs.\ \textit{ilium}), abbreviations are expanded (e.g.\ BP $\rightarrow$ blood pressure), and drug and anatomical names are corrected to standard spellings. Second, speaker role assignment: generic labels \texttt{SPEAKER\_00}/\texttt{SPEAKER\_01} are mapped to \textsc{physician} or \textsc{patient} based on full-transcript context, exploiting clinical questioning patterns, examination language, and symptom-reporting cues. The output schema is \texttt{\_LLMTranscript}, containing a list of \texttt{\_LLMSegment} (fields: index, speaker\_role, cleaned\_text).

Combining role assignment with transcript correction in a single LLM pass avoids a redundant API call and improves accuracy through access to global dialogue context \cite{efstathiadis2024diarization, adedeji2024sound}. The original ASR text is preserved in \texttt{PostProcessedSegment\discretionary{}{}{}.original\_text} for auditability and downstream error analysis.

\subsubsection{Clinical Extraction and FHIR Resource Generation (Step 5)}

Step 5 accepts the post-processed, role-labeled transcript and uses GPT-5-mini via agno Agent with a Pydantic extraction schema, followed by programmatic FHIR construction via \texttt{fhir.resources}, to produce a validated FHIR R4 Bundle. A single LLM pass jointly performs: named entity recognition across four resource types (Condition, MedicationStatement, Observation, Procedure); negation detection (e.g.\ ``no chest pain'' $\rightarrow$ \texttt{verificationStatus=refuted}; ``possible X'' $\rightarrow$ \texttt{unconfirmed}); relation extraction to capture drug–dose, route, and frequency as structured dosage fields; implicit SOAP classification (symptoms $\rightarrow$ unconfirmed Condition; examination/laboratory findings $\rightarrow$ Observation; diagnoses $\rightarrow$ confirmed Condition; plan items $\rightarrow$ MedicationStatement/Procedure); and best-effort terminology coding using SNOMED CT for conditions and procedures, RxNorm for medications, and LOINC for observations.

The LLM output is constrained to the following Pydantic schema:
\begin{itemize}
  \item \texttt{\_Condition}: text, clinical\_status, verification\_status, snomed\_code, icd10\_code, segment\_indices
  \item \texttt{\_Medication}: drug\_name, dose, route, frequency, status, rxnorm\_code, segment\_indices
  \item \texttt{\_Observation}: text, value, unit, loinc\_code, segment\_indices
  \item \texttt{\_Procedure}: text, status, snomed\_code, segment\_indices
\end{itemize}

FHIR R4 resources are then constructed programmatically via \texttt{fhir.resources} Pydantic models rather than as free-form JSON. The bundle includes Encounter, Condition, MedicationStatement, Observation, and Procedure resources, assembled as a collection-type Bundle. The Encounter period is derived from the minimum and maximum timestamps of transcript segments. Each resource carries a custom extension encoding \texttt{segment\_indices}, linking the resource back to its source transcript segments for auditability.

This unified approach trades an accuracy ceiling (reported at 91\% semantic completeness and 0.88 interoperability for dedicated transformer-based NER pipelines with ontology-grounded normalization \cite{brens2025semantic}) for engineering simplicity and end-to-end pipeline integration. This trade-off is consistent with the demonstrated feasibility of LLM-driven FHIR transformation in the literature \cite{hong2020nlp2fhir, li2024fhirgpt}.

\subsubsection{FHIR Validation and Provenance Tracking (Step 6)}

Step 6 performs structural validation of the generated bundle and attaches provenance metadata to each clinical resource, using \texttt{fhir.resources} Pydantic R4B models \cite{hl7fhir2022r4b}. Validation parses each resource through its corresponding \texttt{fhir.resources} class (Condition, MedicationStatement, Observation, Procedure, Encounter, Bundle). Pydantic \texttt{ValidationError} exceptions are caught and converted to \texttt{ValidationIssue} records (fields: severity, resource\_type, resource\_id, message). The bundle is flagged \texttt{valid=True} if and only if no error-severity issues are raised; warning-severity issues are recorded but do not invalidate the bundle.

For each clinical resource, a FHIR \texttt{Provenance} resource is generated and appended to the bundle, carrying: a \texttt{target} reference to the clinical resource; a \texttt{recorded} timestamp of generation; an \texttt{agent.who} identifying the pipeline component (\path{rtm-pipeline/step05/gpt-5-mini}); an \texttt{entity.what} referencing the source audio filename; and a custom \texttt{nlp-system} extension valueString encoding the pipeline version. This approach follows the provenance extension conventions established by the SMART Text2FHIR pipeline \cite{miller2023smart}, enabling downstream consumers to reconstruct NLP provenance without rerunning the pipeline \cite{frei2025inferno}.

\subsection{Implementation Details}

The pipeline is implemented in Python 3.13, managed with uv. All inference is CPU-only: pyannote.audio and Whisper run on \texttt{device="cpu"} with \texttt{torch.float32}, with the ASR model requiring approximately 6\,GB of RAM for the 1.55B-parameter checkpoint. LLM calls are made to the OpenAI API (configured via the \texttt{OPENAI\_API\_KEY} environment variable); diarization model weights are retrieved from the HuggingFace Hub (configured via \texttt{HF\_TOKEN}). The runtime environment is NixOS, managed through devenv with direnv-based activation.
